\subsection{Approximation of continuous functions by step functions}

PROPOSITION 19. --- Let X be a locally compact space, \(\Phi\) a clan of subsets of X, containing the set of compact subsets of X. For every continuous mapping f of X into a Banach space F (resp. every continuous, real-valued function \(f \geqslant 0\) on X) with compact support K, there exists a sequence \((\mathbf{g}_n)\) of functions in \(\mathcal{E}_{\mathrm{F}}(\Phi)\) with support contained in K (resp. a sequence \((g_n)\) of functions in \(\mathcal{E}(\Phi)\) such that \(0 \leqslant g_n \leqslant f\) for every n) that converges uniformly to f (resp. \(f\) ).

Since \(\mathbf{f}\) is uniformly continuous on \(K\) , one can cover \(K\) by a finite number of compact sets \(\mathbf{M}_i\) ( \(1 \leqslant i \leqslant m\) ) such that the oscillation of \(\mathbf{f}\) on each \(\mathbf{M}_i\) is \(\leqslant 1/n\) . Since the \(\mathbf{M}_i\) and \(K\) belong to \(\Phi\) , there exists a partition of \(K\) into sets \(\mathbf{N}_j \in \Phi\) such that each of the sets \(\mathbf{M}_i \cap K\) is the union of a certain number of the \(\mathbf{N}_j\) (No.~9, Lemma). Let \(\mathbf{a}_j\) be an element of \(F\) such that \(|\mathbf{f}(x) - \mathbf{a}_j| \leqslant 1/n\) on \(\mathbf{N}_j\) . Setting \(\mathbf{g}_n = \sum_j \mathbf{a}_j \varphi_{\mathbf{N}_j}\) , we have \(|\mathbf{f} - \mathbf{g}_n| \leqslant 1/n\) , whence the proposition in this case. One argues similarly for a continuous real-valued function \(f\) , on taking \(a_j = \inf_{x \in \mathbb{N}_j} f(x)\) and \(g_n = \sum_j a_j \varphi_{\mathbf{N}_j}\) .

COROLLARY 1. --- Let \(\mu\) be a measure on X; the space \(\mathcal{E}_{\mathrm{F}}(\Phi)\) is dense in each of the spaces \(\mathcal{L}_{\mathrm{F}}^{p}\) ( \(1 \leqslant p < +\infty\) ).

For, it follows from Prop. 19 and the criterion for convergence in mean for uniform limits of functions with compact support (§3, No.~3, Prop. 4) that \(\mathcal{E}_{\mathrm{F}}(\Phi)\) is dense, for the topology of convergence in mean of order \(p\) , in the closure of the space \(\mathcal{K}_{\mathrm{F}}\) of continuous functions with compact support, whence the corollary.

COROLLARY 2. --- For every closed subset S of X, every function \(f \in \mathcal{K}(X, S; \mathbf{C})\) is the uniform limit of linear combinations \(\sum_{i} \lambda_i \varphi_{K_i}\) , where the \(\lambda_i\) belong to \(\mathbf{C}\) and the \(K_i\) are compact subsets of \(S\) .

The set A of such linear combinations is a C-algebra. Let \(\Phi\) be the set of subsets M of X such that \(\varphi_{M} \in A\) ; \(\Phi\) is thus a clan all of whose elements are subsets of S, containing the compact subsets of S, and \(\mathcal{E}_{\mathbf{C}}(\Phi) \subset \mathcal{A}\) . It then suffices to apply Prop. 19 to the locally compact space S and the clan \(\Phi\) .

COROLLARY 3. --- If \(\mu\) and \(\nu\) are two measures on X such that \(\mu(K) = \nu(K)\) for every compact subset \(K\) of \(X\) , then \(\mu = \nu\) .

For, it follows from Cor. 2 and the definition of a measure that, for every compact subset S of X, \(\mu\) and \(\nu\) take on the same values in \(\mathcal{K}(\mathrm{X},\mathrm{S};\mathbf{C})\) .

